\documentclass{article}
\usepackage{decomp}
\newenvironment{reflist}{\par\small\setlength{\parindent}{0pt}\everypar{\hangindent=1.6em\hangafter=1}}{\par}
\begin{document}
\decompnote{Collection markers: a cited work that is in the local collection is tagged \localcopy{<key>} after its reference entry; it is decompiled in decompiled/<key>/. A reference to another version of a local work (e.g. its preprint) is tagged \localother{<key>}. A reference without that tag is not in the local collection and needs an external pull. Tags added by \texttt{\detokenize{tools/mark_local.py}}.}

\sourcefile{hua2015.pdf}
\decompnote{Publisher version of record (Nature Communications typeset PDF, 7 pages). The running header (``ARTICLE'' / ``NATURE COMMUNICATIONS \textbar{} DOI: 10.1038/ncomms8923''), the running footer (journal citation line and page number) and the ``© 2015 Macmillan Publishers Limited. All rights reserved.'' line at the foot of every page, and the Nature Communications logo, are left out of the page flow; the footer line and the copyright line are given once below. The Creative Commons licence statement is on p.\,7. The PDF's text layer uses Symbol-font glyphs that it reads as other letters (µ as ``m'', ∼ as ``B'', $<$ as ``o'', $>$ as ``4'', × and − as blanks, = as ``¼'', Θ as ``Y'', ° missing); they are transcribed here from the page images.}

\pdfpage{1}
\noindent ARTICLE\par
\noindent Received 6 Apr 2015 \textbar{} Accepted 25 Jun 2015 \textbar{} Published 3 Aug 2015\qquad DOI: 10.1038/ncomms8923\qquad OPEN\par
\noindent NATURE COMMUNICATIONS \textbar{} 6:7923 \textbar{} DOI: 10.1038/ncomms8923 \textbar{} www.nature.com/naturecommunications\par
\noindent © 2015 Macmillan Publishers Limited. All rights reserved.\par

\section*{Large-volume \emph{en-bloc} staining for electron microscopy-based connectomics}

\noindent Yunfeng Hua\textsuperscript{1}, Philip Laserstein\textsuperscript{1} \& Moritz Helmstaedter\textsuperscript{1}

\begin{quote}\small
\noindent\textsuperscript{1}Department of Connectomics, Max Planck Institute of Brain Research, Max-von-Laue-Strasse 4, D-60438 Frankfurt, Germany. Correspondence and requests for materials should be addressed to Y.H. (email: yunfeng.hua@brain.mpg.de) or to M.H. (email: mh@brain.mpg.de).
\end{quote}
\decompnote{The affiliation and correspondence footnote is printed at the foot of p.\,1; it is placed here after the author list.}

\begin{quote}
\noindent Large-scale connectomics requires dense staining of neuronal tissue blocks for electron microscopy (EM). Here we report a large-volume dense \emph{en-bloc} EM staining protocol that overcomes the staining gradients, which so far substantially limited the reconstructable volumes in three-dimensional (3D) EM. Our protocol provides densely reconstructable tissue blocks from mouse neocortex sized at least 1\,mm in diameter. By relaxing the constraints on precise topographic sample targeting, it makes the correlated functional and structural analysis of neuronal circuits realistic.
\end{quote}

\pdfpage{2}
\noindent Mapping neuronal circuits densely is the goal of connectomics. Current imaging methods in high-resolution connectomics require \emph{en-bloc} staining of neuronal tissue for large-volume EM\textsuperscript{1}. Establishing high-contrast staining in every neurite and synapse in large volumes homogeneously is, however, a substantial challenge to the employed staining protocols. In spite of research into the chemistry of heavy metal deposition in biological samples over several decades, the limitation of the employed protocols to small-dimension samples (typically ∼100–200-µm maximum penetration depth) had been widely accepted. Only very recently, a protocol aimed at whole-brain staining has been proposed\textsuperscript{2}, which, however, is focused on very large whole-brain-sized samples and involves staining procedures in the scale of months, not days. We developed a staining protocol that overcomes the limitations of sample size with a simple, fast and reliable process, which provides high-contrast full-neuropil staining for large (millimetre size) volumes.
\decompnote{The first word is set with a large drop cap ``M''.}

\subsection*{Results}
The widely employed existing EM staining protocols either are limited to small volumes, yielding prohibitive staining gradients for larger structures (Fig. 1a-c), or stain only a subset of the neurites (myelinated fibre staining, Fig. 1d), with the additional requirement for charging compensation\textsuperscript{2}. Our protocol provides, however, high-contrast full-neuropil staining for large (millimetre size) volumes (Fig. 1e). Its chemical logic, key steps and usability for dense circuit reconstruction are described in the following.

Neuronal tissue staining for high-resolution connectomics requires the deposition of heavy metal compounds into the membranes that outline neuronal cell bodies, their dendrites, axons and synaptic vesicles. Osmium (Os) and uranium compounds have been the most widely employed membrane contrasting agents for decades\textsuperscript{4–6}, and uranyl acetate (UA) has also been implied in labelling proteins such as those in the postsynaptic density\textsuperscript{5}. Alterations to both the Os tetroxide (\ce{OsO4}) and UA impregnation steps are at the core of our protocol.

To achieve high-contrast staining throughout large tissue blocks, the relevant compounds have to diffuse from the exterior through many successive lipid membranes to the centre of the tissue block. Secondly, these compounds have to react with the target structures in the membranes (or be deposited in the membrane by other means). There is a tradeoff between these two goals, since processes that enhance membrane attachment will reduce the availability of compounds for further diffusion, and \emph{vice versa}. This is especially true when enhancing reactions are employed that further increase compound attachment, but thus decrease diffusability. One important enhancing reaction is the application of ferrocyanide for enhanced Os staining, which will be discussed below.

In brief, our protocol alterations therefore centre around improving the optimum between tissue penetration and membrane contrast. By separating the \ce{OsO4} penetration from the ferrocyanide enhancement reaction, by leaving out the intermittent washing step and by using a two-step temperature protocol in the UA impregnation, we achieve our homogeneous high-contrast penetration results as shown in Fig. 1e. Most notably, our protocol completely abolishes the intense staining ring generated in the state-of-the-art \emph{en-bloc} staining protocols (Fig. 1b,c) which had been implied in blocking further stain penetration. When imaged at high resolution (Fig. 2), we find that samples stained with our protocol show intense and continuous membrane contrast (Fig. 2e,f and Supplementary Fig. 1) at about 500\,µm from the sample surface (Fig. 2a) with no indication of contrast decrease over this penetration distance.

Our protocol alterations are based on three notions about the underlying diffusion and reaction steps. We report these notions here to illustrate the reasoning that guided our protocol development (Fig. 3). First, we assumed that the classic \ce{OsO4}-alkene reaction mechanism (the Sharpless dihydroxylation reaction\textsuperscript{7–9}) is at the core of the \ce{OsO4}–membrane interactions (Fig. 3a). Second, we find indications that in addition to the covalent binding of \ce{OsO4} to unsaturated lipids, quadruply reduced Os\textsuperscript{IV}, made available as \ce{OsO2} by dismutation of Os\textsuperscript{VI} compounds (Fig. 3b), and dissolved in the membrane, substantially contributes to membrane contrast. Finally, we exploit the fact that the diffusability and reactivity of UA is differentially affected by temperature, such that lower temperature reduces reactivity\textsuperscript{10} more strongly than it reduces diffusibility.

These reaction mechanisms imply that \ce{OsO4} diffuses across membranes by switching between its ionized hydrophilic form \ce{[OsO4(OH)2]^{2-}} and its non-polar lipophilic form \ce{OsO4} (Fig. 3a, step 1 when entering the membrane, inverted step 1 (`−1') when exiting the membrane, see sketch in Fig. 3c). After entering the membrane, \ce{OsO4} can form cyclic osmate esters by covalent binding to unsaturated membrane lipids\textsuperscript{11–13} (Fig. 3a, step 2). However, these intermediates slowly undergo dihydroxylation\textsuperscript{11,14} (Fig. 3a, step 3), with the effect of Os leaving the membrane again, yielding \ce{[OsO2(OH)4]^{2-}}, a water soluble Os compound at oxidation level VI. This step alone would result in membrane de-staining; however, it has been suggested based on chemical analysis\textsuperscript{12} that instable Os\textsuperscript{VI} compounds dismutate quickly to \ce{OsO4} and \ce{OsO2} (Fig. 3b, step 5). \ce{OsO2} is non-polar and thus likely lipophilic and will be dissolved in the membrane, thus securing in-membrane Os deposition.

Following this logic, the main effect of reducing agents such as potassium ferrocyanide (\ce{K4Fe(CN)6}, step 4 in Fig. 3a) is to convert the VIII-oxidized water-soluble Os form into a VI-oxidized water-soluble form\textsuperscript{15}, which makes more Os\textsuperscript{VI} available for the dismutation step (Fig. 3b, step 5), and thus generates additional \ce{OsO2} to be deposited in the membrane. This pathway circumvents the osmate ester formation step (Fig. 3a, step 2), which requires unsaturated membrane lipids as reaction partners\textsuperscript{15}.

Together, this would imply that the availability of Os\textsuperscript{VIII} for sequential membrane uptake/loss events (steps 1, −1) is essential for the penetration of membrane stain into the tissue blocks. The generation of reduced Os species (Os\textsuperscript{VI} and Os\textsuperscript{IV}) only via the dihydroxilation (Fig. 3a), however, may saturate with decreasing availability of unsaturated membrane lipids. Both effects could explain why the OTO protocol\textsuperscript{16} (which involves two \ce{OsO4} incubation steps and an intermittent thiocarbohydrazide (TCH) enhancing reaction, but no reducing agent, Fig. 3c) provides good stain penetration but low contrast (limited by availability of unsaturated membrane lipids). In addition, \ce{OsO4} can be washed out in this protocol during the washing steps following \ce{OsO4} incubation (via reaction −1).

On the other hand, reduced Os protocols (rOTO\textsuperscript{17}; using ferrocyanide/\ce{OsO4} mixture instead of \ce{OsO4}, also used in recent studies\textsuperscript{18,19}, Fig. 3d) `push' almost the entire water-dissolved Os\textsuperscript{VIII} fraction into the membrane thus avoiding unsaturated lipid depletion and increasing local contrast, but hampering further membrane penetration. This can explain why the reduced Os staining protocols (Fig. 3d) provide very good contrast in the periphery, but with limited penetration (typically $<$200\,µm from the tissue surface).

This logic motivates our alteration of the staining protocol (Fig. 3e): we first allow \ce{OsO4} to penetrate deep into the tissue (via steps 1, -1) and then immediately apply the reducing agent ferrocyanide (without an intermediate washing step, which would extract Os\textsuperscript{VI} or Os\textsuperscript{VIII} from the membrane).

\pdfpage{3}
\begin{center}
\includegraphics[width=0.95\linewidth,height=0.55\textheight,keepaspectratio]{figures/fig1.png}
\end{center}
\noindent\textbf{Figure 1 \textbar{} Novel large-scale \emph{en-bloc} EM staining protocol for dense connectomic circuit reconstruction: protocol comparison.} (\textbf{a}) Screening strategy for brain biopsy samples $>$1\,mm in smallest dimension. Samples were screened after cutting the sample in about half, such that the core of the sample was exposed and could be tested for staining quality. Samples were screened in overview (1), periphery (2), intermediate (3) and core regions (4). (\textbf{b–e}) Scanning electron microscopy images from mouse cortex biopsy samples screened as indicated in \textbf{a} (regions 1–4, left to right columns) and stained with \emph{en-bloc} protocols from ref. 18 (\textbf{b}, rOTO), ref. 19 (\textbf{c}, rOTO) and ref. 2 (\textbf{d}, PATCO) and this protocol (\textbf{e}) Note the good staining quality in the periphery (first column) in \textbf{b,c} but the strong over-staining in the intermediate region yielding stain uptake of a subset of neurons (\textbf{b,c}) and a substantial staining gradient, which is relieved in this protocol. (\textbf{e}) PATCO provides good stain penetration but stains myelinated neurites only, (\textbf{d}) Scale bars, 200\,µm in \textbf{b–e} left column and 2\,µm in other columns.
\decompnote{As printed, the panel letters at the end of the caption look misplaced: the sentence ``PATCO provides good stain penetration but stains myelinated neurites only'' describes panel d (PATCO, ref. 2), but is preceded by ``(e)'' and followed by ``(d)''. The caption also calls the first column the periphery, while (a) and ``regions 1–4, left to right columns'' make the first column the overview (1) and the second the periphery (2). Kept as printed.}
\alttext{Figure 1 fills the upper two thirds of p.\,3. (a) Schematic, one row, left to right with black arrows between the steps: a white upright cylinder labelled ``Brain biopsy'' with a dimension arrow ``1 mm'' across its top (diameter) and a vertical dimension arrow ``$>$1 mm'' along its height; a dark grey cylinder labelled ``EM staining''; ``Blockface SEM'': a short grey cylinder with a white top face and, above it, a tilted dark grey piece (the cut-off half); a large light grey disc (the exposed sample face) with a black square outline over its upper left part, the number ``1'' at the square's left side, and three small squares inside the square, labelled ``2'', ``3'' and ``4'', running diagonally from upper left (near the disc edge) to lower right. (b–e) Four rows of four square greyscale scanning electron micrographs, one row per protocol (rows labelled b, c, d, e at upper left); the first column is a low-magnification overview of the sample corner (sample edge as a curved boundary at upper left), columns 2–4 are high-magnification views of neuropil. Scale bars only in row b: a white bar at the bottom of the overview (200 µm per caption) and a white bar at the bottom right of the second image (2 µm per caption). Row b: the overview shows a broad dark band following the curved sample edge a short distance inside it, with lighter grey tissue inside the band and further in; column 2 has crisp dark membranes and a myelinated axon at the upper left; column 3 is darker and blotchy with poorly defined membranes and a large very dark elongated profile at the lower left; column 4 is mid-grey with lower membrane contrast and a few dark myelin profiles. Row c: the overview shows a black rim and dark band at the edge and grey tissue inside with a faint darker zone; column 2 has high contrast with several thick dark myelin rings; column 3 has moderate contrast and a very dark elongated cell profile in the centre right; column 4 is pale grey with low membrane contrast. Row d: the overview is light grey and nearly uniform (bright background outside the edge at upper left); columns 2–4 have a pale grey, low-contrast background in which only dark ring-shaped myelinated axon profiles stand out (more of them in columns 3 and 4). Row e: the overview shows a dark rim at the edge and homogeneous mid-grey tissue with many lighter round spots further in; columns 2–4 all show uniformly high-contrast neuropil with dark continuous membranes, a thick black myelinated fibre bundle running diagonally through column 2 and dark myelin rings in columns 3 and 4.}

\noindent How did we arrive at this protocol alteration? Initially, we attempted to use periodic acid (which was so far used to stain myelinated fibres\textsuperscript{2,20}), however, together with \ce{OsO4}, such that periodic acid could attach to dihydroxylated lipids and make the binding site again available for a further round of Os amplification through TCH. However, the effect of periodic acid application after \ce{OsO4} incubation (Fig. 3f) was opposite: membrane staining contrast got worse (Fig. 3g). We also noted that periodic acid weakened the strong macroscopic tissue staining (black under bright-field illumination, Fig. 3f). Since the most likely source of this heavy black staining is membrane-dissolved or -precipitated \ce{OsO2} (`Os black'), we therefore wondered whether in fact periodic acid had oxidized \ce{OsO2} to water soluble and colourless \ce{[OsO4(OH)2]^{2-}}.

This made us consider an addition to the classic concept of Os membrane staining. So far, specific chemical binding of \ce{OsO4} to unsaturated lipids was considered the main source of long-term membrane staining\textsuperscript{11,12}. But what if membrane-dissolved black-colored \ce{OsO2} was a parallel and crucial source of membrane contrast? Thus the macroscopically black appearance of the \ce{OsO4}-impregnated brain (Fig. 3f, middle) could result from \ce{OsO2} enriched in membranes. A simple water/octanol partition test (Fig. 3h) provided support for this hypothesis: when mixing equal parts of water-dissolved 2\% \ce{OsO4} and octanol, the organic and aqueous phases separated (Fig. 3h, left). After several hours, however, the organic phase had turned deep black (Fig. 3h, right). Since 1-octanol can progressively reduce \ce{OsO4} to \ce{OsO2}, but cannot form stable osmate esters due to the absence of carbon–carbon double bonds the partition of black compounds in
\pdfpage{4}
the lipophilic phase as seen in the partition test (Fig. 3h) is likely caused by the strongly lipophilic \ce{OsO2}.

\begin{center}
\includegraphics[width=0.6\linewidth,height=0.55\textheight,keepaspectratio]{figures/fig2.png}
\end{center}
\noindent\textbf{Figure 2 \textbar{} High-resolution screening of ultrastructure and staining contrast at the core of ∼1-mm-sized tissue sample.} (\textbf{a–d}) Position of screening region within sample. (\textbf{e,f}) Two example high-resolution scanning electron microscopy images (imaged at voxel size of 1.35\,×\,1.35\,nm) showing continuous strong membrane contrast in plasma and vesicle membranes, indications of postsynaptic density, and intracellular organelles. Images in \textbf{e,f} were imaged first on a freshly cut block face, and images in \textbf{a–d} were taken afterwards to avoid effects of multiple exposure in the high-resolution images. See Supplementary Fig. 1 for further high-resolution screening results. Scale bars, 200\,µm in \textbf{a}, 40\,µm in \textbf{b}, 2\,µm in \textbf{c,d} and 500\,nm in \textbf{e,f}.
\alttext{Figure 2 is in the upper left column of p.\,4, above its caption. Six greyscale scanning electron micrographs: a narrow left column with a, b (small, side by side at top), c and d (larger, stacked), and a right column with e (top) and f (bottom). (a) Low-magnification view of the block face: a light grey square, slightly rotated, on a dark background; a white square outline labelled ``b'' near the centre; white scale bar at lower right (200 µm per caption). (b) Enlargement of that region: grey tissue with many lighter round spots and a dark dot at lower left; two small white squares labelled ``c'' and ``d'' at the top; white scale bar at lower right (40 µm). (c, d) Neuropil at medium magnification: many profiles outlined by dark, continuous membranes, dark mitochondria, myelinated axons as dark rings; each has a white square outline near the centre and a black scale bar at lower right (2 µm). (e, f) High-magnification images: neurite profiles bounded by dark, sharp, continuous membranes; clusters of small round vesicle profiles fill some profiles; a mitochondrion with visible cristae in the centre of e and two in f (upper left and bottom); in f a stack of dark parallel membranes at the right; white scale bars at lower right (500 nm).}

In summary, with respect to Os staining, our protocol follows the logic that \ce{OsO4} needs to stay in VIII-oxidized form for subsequent membrane uptake events and thus tissue penetration; but it needs to be reduced to oxidation level IV (\ce{OsO2}) to be kept in the membrane for high-contrast tissue staining.

With respect to uranium staining, we hypothesized that the temperature dependency of UA–substrate interactions was stronger for the chemical reactions leading to binding of UA to protein and membrane targets\textsuperscript{10} than for its diffusibility through tissue. We therefore exposed the samples to a two-step UA immersion: first overnight at low temperature (4\,°C) to allow diffusion over large distances; and then a brief step to 50\,°C (for 2\,h) to enhance reactivity with protein and membrane targets. The latter enhanced staining contrast (Supplementary Fig. 2). In addition, we optimized the dehydration and resin-embedding steps for larger tissue volumes (see Methods).

The protocol was then controlled for inter-experimentator variability and for variability between cortical areas (four independent experimentators; somatosensory and parietal cortex of mouse, Fig. 1e, Supplementary Fig. 3, Supplementary Table 2).

We finally tested whether the new protocol in fact provides sufficient staining contrast for dense circuit reconstruction (Fig. 4). For this, we acquired a three-dimensional (3D) serial block-face electron microscopy (SBEM) image data set sized 65\,×\,51\,×\,41\,µm\textsuperscript{3} with a voxel size of 12\,×\,12\,×\,30\,nm\textsuperscript{3} (Fig. 4b, Supplementary Movie) from the core of one 1-mm-sized sample (at about 500\,µm from the sample border, Fig. 4a) and asked 7 trained annotators to start at 116 random neurite seed points in the data set (see Methods for details of the sampling procedure). From each seed point, the annotators were asked to reconstruct the respective neurite within a cubic bounding box of 10\,µm edge length (Fig. 4b,c; average traced path length per seed was 27.9\,µm including branches), yielding a total tracing of 22.7\,mm path length. We then applied the RESCOP skeleton consensus computation\textsuperscript{21}, which measures the distribution of agreement between tracers (vote histogram, Fig. 4d). We then used RESCOP to determine the traceability prior (Fig. 4f), which reports the distribution of tracing difficulty in the data set, by fitting to the measured vote histogram (Fig. 4d reports the measured and Fig. 4e the fitted vote histogram). Finally, RESCOP provides a prediction of tracing accuracy in dependence of the number of redundant annotations (Fig. 4g). This quantification allows the comparison to tracing precision in both the published retina data sets\textsuperscript{18,21,22} and to a data set acquired from mouse somatosensory cortex but stained with the conventional \emph{en-bloc} protocol (protocol from ref. 18, see Fig. 1b; data set: Boergens \emph{et al.}, unpublished data set). We find that the dense reconstructions performed at the core of a sample stained with our new protocol are at least as good as those from the retina data set and a conventionally-stained cortex data set (Fig. 4g). Thus, our staining quality is comparable to the previous (volume limited) protocols and allows for dense neuronal circuit reconstruction.

\subsection*{Discussion}
The reported EM staining protocol is currently mainly aimed at yielding consistently high-contrast staining for samples sized about a millimetre in diameter. This is an enormous alleviation of the experimental burden in EM-based connectomics. In many experimental settings, the neuronal tissue-of-interest is a structure of several hundred micrometres in diameter (for example, barrels in mouse S1 cortex, subnuclei in the thalamus or directional preference columns in visual cortex). With the existing protocols, the maximum sample size was just on that same spatial scale, that is, one had to extract samples of about 400\,µm in diameter at about 10–50\,µm targeting precision from entire mouse brains (Fig. 4h). This represented a substantial obstacle in \emph{en-bloc} EM imaging. Our protocol enables sample volumes that considerably relieve the precision requirements in sample targeting (Fig. 4h).

While such obstacles may be just acceptable for purely structural investigations (where multiple samples can be readily obtained at low resource consumption), this sampling statistics was close to prohibitive for correlated functional/structural experiments. In these settings, the functional experiment alone has usually challenging success statistics (often one-fifth successful experiments). Increasing this sampling burden by a factor of about 5–10 makes such correlated experiments close to impossible. Therefore, our protocol not only makes large-scale structural imaging possible, but substantially alleviates the EM \emph{en-bloc} staining challenges encountered in correlated experimental settings, opening the path to many promising structure-function studies in the emerging field of connectomics.

We tested our protocol for samples sized about 1\,mm in diameter. In principle, an extension of the protocol to samples at the size of a hemisphere or an entire mouse brain should be possible by prolonging the respective impregnation steps. With larger samples, however, macroscopic sample damage becomes more likely, such as cracks in the tissue blocks. With high-contrast staining, another challenge is maintaining conductivity
\pdfpage{5}
throughout large sample volumes, since the less metal deposition in the cytosol (less `background') the less conductive the sample. Using conductive resins or increasing cytosol staining may relieve these challenges for very large samples.

\begin{center}
\includegraphics[width=0.95\linewidth,height=0.55\textheight,keepaspectratio]{figures/fig3.png}
\end{center}
\noindent\textbf{Figure 3 \textbar{} Possible chemical reactions underlying Os-based membrane staining and logic of our staining protocol.} (\textbf{a}) Sharpless-dihydroxilation (modified from ref. 9) and possible effect of reducing agents such as ferrocyanide. Note that reaction steps 1–5 are referred to in Results, and inverse reaction steps are referred to as `−1' to `−5'. (\textbf{b}) Os\textsuperscript{IV} dismutation suggested to be critical for deposition of \ce{OsO2} in the membrane (after ref. 12) (\textbf{c–e}) Possible reaction and diffusion steps involved in stain penetration and contrast enhancement in the conventional OTO protocol, (\textbf{c}) rOTO protocol (\textbf{d}, see also Fig. 1b,c) and this protocol (\textbf{e}, compare with Fig. 1e). Two membranes are visualized to indicate penetration of successive plasma and intracellular membrane bilayers. Only key reaction steps are indicated. (\textbf{f}) Simple test experiment that lead us to consider \ce{OsO2} as a main source of membrane contrast. After Os impregnation, the tissue obtains a dark black colour under bright-field illumination (middle). This is reversed by periodic acid application (right), suggesting that \ce{OsO2} is oxidized to colourless Os species. (\textbf{g}) EM contrast after periodic acid application was decreased, staining only myelinated fibres (inset, note region between myelinated fibres has very low contrast making the detection of single neurites impossible), suggesting that \ce{OsO2} had been essential for Os membrane contrast. (\textbf{h}) Simple octanol/water partition test: equal fractions of \ce{OsO4} in aqueous solution and octanol partition into organic (top) and aqueous (bottom) phase. After several hours (right) the lipophilic phase is deep black, supporting the notion that \ce{OsO2} (which is the product of \ce{OsO4} reduction by octanol) can be dissolved in lipophilic membranes and may contribute crucially to membrane contrast. Scale bars, 20\,µm and 5\,µm in \textbf{g} and insert, respectively.
\decompnote{The caption says ``Os\textsuperscript{IV} dismutation'' for panel b, but panel b and the main text (``dismutation of Os\textsuperscript{VI} compounds (Fig. 3b)'') show the dismutation of Os\textsuperscript{VI} (2 Os\textsuperscript{VI} → Os\textsuperscript{IV} + Os\textsuperscript{VIII}). Kept as printed.}
\alttext{Figure 3 fills the upper three fifths of p.\,5. Panels a and b (top left and top centre) are chemical reaction schemes; c, d, e (lower row) are cartoons; f, g, h (top right) are photographs and a micrograph.
(a) Reaction scheme with a horizontal dashed line across a and b separating ``Organic'' (above) from ``aqueous'' (below); circled numbers mark the steps. Step 1 (on the dashed line): the dianion \ce{[OsO4(OH)2]^{2-}} drawn below the line (Os in the centre, two axial =O, four OH; labelled VIII) gives \ce{OsO4} above the line (tetrahedral Os with four =O, labelled VIII), releasing ``2 OH$^-$''. Step 2: \ce{OsO4} reacts with an alkene R–CH=CH–R (drawn top centre) to the cyclic osmate ester at top left (a five-membered Os(=O)\textsubscript{2}–O–C–C–O ring, Os labelled VI, the two R groups on wedge/hash bonds). Step 3 (on the dashed line, with ``2 OH$^-$'' and ``2 H\textsubscript{2}O'' entering from the aqueous side): the ester gives the diol HO–CH(R)–CH(R)–OH (above the line) and the Os(VI) dianion \ce{[OsO2(OH)4]^{2-}} below the line (Os with two axial =O and four OH, labelled VI). Step 4 (below, between the two dianions): ``2K\textsubscript{4}Fe(CN)\textsubscript{6}'' → ``2K\textsubscript{3}Fe(CN)\textsubscript{6}'', with an arrow from the Os(VIII) dianion to the Os(VI) dianion. (b) Step 5 (on the dashed line): ``2'' Os(VI) dianions (labelled VI) react to give O=Os=O (labelled IV) above the line and one Os(VIII) dianion (labelled VIII) below it, ``+ 2 OH$^-$ 2 H\textsubscript{2}O''.
(c) Left label ``OsO\textsubscript{4} staining''. Three stacked stages (separated by large grey downward arrows) of two lipid bilayers (drawn as two rows of grey head-group circles with wavy tails), left and right, with a bottom arrow ``Penetration depth'' pointing right. Stage 1, ``OsO\textsubscript{4}'': open yellow circles (\ce{[OsO4(OH)2]^{2-}}) outside the membranes; red arrows ``1'' (into the left membrane, where filled yellow circles = \ce{OsO4} are), ``−1'' (out of it towards the right) and ``1'' (into the right membrane); a curved red arrow ``2'' inside the left membrane leads to open black circles (\ce{OsO3}-R). Stage 2, ``Wash'': red arrows ``−1'' pointing out of both membranes to open yellow circles; a curved arrow ``3+5'' in the left membrane leads to a filled black circle (\ce{OsO2}); both membranes contain some black dots, open black circles and filled yellow circles. Stage 3, ``After wash'': only a few black dots (\ce{OsO2}) remain, three in the left and two in the right membrane.
(d) Left label ``Reduced Os staining''. Stage 1, ``OsO\textsubscript{4} K\textsubscript{4}Fe(CN)\textsubscript{6}'': open yellow and green circles (\ce{K4Fe(CN)6}) mixed outside the left membrane; two curved red arrows labelled ``4+5'' end on black dots at the left membrane surface; the left membrane holds four black dots and two yellow dots, the right membrane only one yellow dot. Stage 2, ``After wash'': six black dots in the left membrane, two in the right; arrow ``Penetration depth''. Legend below d: filled yellow circle ``OsO\textsubscript{4}'', open yellow circle ``[OsO\textsubscript{4}(OH)\textsubscript{2}]\textsuperscript{2–}'', filled green circle ``K\textsubscript{4}Fe(CN)\textsubscript{6}'', open black circle ``OsO\textsubscript{3}-R'', filled black circle ``OsO\textsubscript{2}''.
(e) Left label ``This protocol''. Stage 1, ``OsO\textsubscript{4}'': as stage 1 of c (arrows ``1'', ``−1'', ``1'' and ``2''). Stage 2, ``K\textsubscript{4}Fe(CN)\textsubscript{6}'': green circles outside and between the membranes; red curved arrows ``−1+4+5'' (from yellow dots in the left membrane via a green dot back to black dots) and an unlabelled one at the right membrane, ``4+5'' at the lower left and ``3+5'' in the left membrane, all ending on black dots; both membranes now contain several black dots. Stage 3, ``After wash'': many black dots in both membranes (seven in the left, five in the right), more than in c and than in the right membrane of d; arrow ``Penetration depth''.
(f) Photographs of three small glass vials, left to right, joined by arrows labelled ``2\% OsO\textsubscript{4} (1 h)'' and ``1\% IO\textsubscript{4}$^-$ (1 h)'': a pale tan tissue piece at the bottom of the first vial (handwritten ``CaC'' on the glass), a black tissue piece in the second (handwritten label, probably ``OsO4''), and a brownish-olive piece in the third (handwritten ``IO\textsubscript{4}$^-$'').
(g) Greyscale SEM micrograph: grey neuropil with many small dark-rimmed round myelinated axon profiles, often in bands, and large pale grey areas between them with little visible structure; a small white square at left centre is enlarged in the inset at lower right (dark rings of myelin around pale axons, with a homogeneous light grey region between them); white scale bars at the bottom of the main image (20 µm) and of the inset (5 µm).
(h) Photograph of two test tubes against a grey background, labelled ``OsO\textsubscript{4}'' (left) and ``OsO\textsubscript{2}'' (right) in white; red dashed lines mark the octanol/water boundary in each tube, labelled ``Octanol'' above and ``water'' below at the left. Left tube: pale yellowish upper phase, clear lower phase. Right tube: black upper (octanol) phase, clear lower phase.}

Astonishingly, our protocol abolishes a decade-long obstacle to EM staining with only a few targeted modifications of \emph{en-bloc} staining. These were, however, based on particular models of the deposition of Os in biological membranes. At the core is the notion that \ce{OsO2} substantially contributes to membrane contrast. This is far from proven, and further experiments may elucidate the contributions of Os at low oxidation states to membrane contrast in biological samples. While only tested for neuropil, our protocol may be well-applicable to other forms of tissue, where 3D EM reconstruction is desirable.

\subsection*{Methods}
\textbf{Animals.} Male mice (P28-P30) were anesthetized and killed by transcardial perfusion of fixative. All procedures were approved by the local animal care and use committee and were in accordance with the laws of animal experimentation issued by the German federal government.

\pdfpage{6}
\begin{center}
\includegraphics[width=0.95\linewidth,height=0.55\textheight,keepaspectratio]{figures/fig4.png}
\end{center}
\noindent\textbf{Figure 4 \textbar{} Dense neurite reconstruction in large tissue blocks.} (\textbf{a,b}) Dense neurite reconstruction in a SBEM data set sized 65\,×\,51\,×\,41\,µm\textsuperscript{3} from the core ($>$500\,µm from the closest sample surface) of a sample from mouse somatosensory cortex. (\textbf{a}) Overview of sample location (stitched from four overview images); (\textbf{b}) 3D display of SBEM-data set boundaries and all locally dense neurite tracings used for tracing test. (\textbf{c}) Two example neurites reconstructed by seven independent annotators. Note complete agreement (right) and local disagreement when tracing spines (left). (\textbf{d,e}) Measured (\textbf{d}) and fitted (\textbf{e}) vote histogram of tracer agreement reporting the total number of tracers (\emph{N}) and the agreeing number of tracers (\emph{T}) for each skeleton-edge (redundancy-corrected, see Methods). (\textbf{f}) fitted distribution of skeleton-edge difficulty (or edge probability) \emph{p}(\emph{p}\textsubscript{e}) from the tracings summarized in \textbf{d}, yielding the fitted vote histogram (\textbf{e}, see RESCOP\textsuperscript{21} and Methods for details). (\textbf{g}) Prediction of tracing accuracy (for full neurite tracings, continuous lines and for focused re-annotation of the disagreeing locations, dashed lines, see ref. 21) from test tracings for this data set compared with published retina tracings (`e1088' from ref. 21) and a comparison cortex data set stained with the protocol as shown in Fig. 1b (Boergens \emph{et al.}, unpublished data set). (\textbf{h}) Illustration of sampling challenge when targeting modules in cortex (here: `barrels' in mouse S1 layer 4, grey cylinder) using a protocol that can only stain about 400-µm-wide samples (black) compared with the procedural relief when obtaining about 1-mm-wide samples (red). Scale bars, 200\,µm (\textbf{a}), 1\,µm (\textbf{c}).
\alttext{Figure 4 fills the upper two fifths of p.\,6. Eight panels: a, b, c, h are images and drawings; d, e, f, g are plots.
(a) Greyscale SEM overview of the block face: a large square-ish light grey tissue area with many small lighter spots and a few dark holes, dark irregular border regions at the left and right edges; a white double-headed arrow labelled ``500 µm'' runs from a small white square labelled ``b'' (left of centre) to the upper right corner of the tissue; white scale bar ``200 µm'' at lower right.
(b) 3D rendering of the SBEM data set as a slightly rotated box with greyscale EM images on its faces; on the front face dense clusters of coloured skeleton tracings (purple, red, green, orange, blue) in several groups; dark round holes (blood vessels) on the faces.
(c) Two coloured skeleton reconstructions drawn vertically, each made of several overlapping thin lines in different colours: the left neurite has several short side branches (spines) along its length where the coloured lines diverge; the right one is a single wavy line along which all colours overlap. An L-shaped black scale bar at lower left (1 µm per caption).
(d, e) Two heat maps side by side, titled ``Meas.'' (d) and ``Fit'' (e), sharing a colour bar at the right labelled ``\# edges'' (log scale, ticks 1, 10, 100, 1,000; jet colours from dark blue at 1 to dark red at the top, about 7,000). x axis ``Total votes for an edge (\emph{N})'', ticks labelled 0, 3, 6; y axis ``Agreeing votes for an edge (\emph{T})'', 0 at top, 3, 6; an 8 × 8 grid of cells (0–7). Cells with T $>$ N (lower left triangle) are dark blue (lowest colour), as is the whole T = 0 row in d. In both maps the diagonal T = N is red/orange (about 600 to 3,000 edges; red at N = T = 1, orange at 2–4, red again at 5–6) and the cell at N = T = 7 is the darkest red (about 6,000 or more); cells above the diagonal are green to yellow (about 25 to a few hundred edges). In e the T = 0 row is filled (yellow at N = 1, cyan-green towards the middle, yellow at N = 7) and the off-diagonal cells near T = 1–2, N = 3–6 are bluer (about 25–50) than in d (about 45–90). Data: fig4d.csv, fig4e.csv.
(f) Line plot. x axis ``\emph{p}\textsubscript{\emph{e}}'', linear, ticks labelled 0 and 1 (and an unlabelled tick at 0.5); y axis ``\emph{p}(\emph{p}\textsubscript{\emph{e}})'', log scale from 10$^{-1}$ to 10$^{3}$ (labels 10$^{-1}$, 10$^{1}$, 10$^{3}$). One blue line that is below the axis almost everywhere and appears as 15 narrow spikes: one at p\textsubscript{e} = 0 reaching about 70, several small ones (about 0.1–2) at p\textsubscript{e} = 0.02–0.14, one at 0.30 (about 0.6), 0.60 (about 2), 0.76 (about 0.9), 0.83 (about 2), 0.91 (about 7), a double spike at 0.94–0.95 (about 3 and 17), small ones at 0.99, and a spike at p\textsubscript{e} = 1 reaching the top of the axis (10$^3$).
(g) Line plot. x axis ``\# annotators'', linear, ticks labelled 1, 10, 20 (the ``1'' tick sits at the left end of the axis); y axis ``Mean inter-error distance (µm)'', log scale 10$^{1}$ to 10$^{7}$ (labels 10$^1$, 10$^3$, 10$^5$, 10$^7$). Legend: red ``e1088 (retina)'', green ``This protocol (cortex)'', blue ``rOTO (cortex)''. Solid lines (1 to 23 annotators) rise slowly, with small wiggles, from about 15–20 µm at 1 annotator to about 235–240 µm at 23 for red and blue and to about 690 µm for green; red is the highest up to 10 annotators, green from 11 annotators on, and blue is the lowest at almost every point. Dashed lines start at 10$^2$ µm (red and green at about 1.7 annotators, blue at about 3.4) and rise steeply and almost straight on the log scale: green reaches 10$^7$ at about 20 annotators, blue at about 22, red bends to a shallower slope above about 10$^5$ and leaves the plot at about 26 annotators near 8 × 10$^6$.
(h) Schematic drawing: a small black-outlined grey cylinder (a barrel) with double-headed arrows ``350'' and ``400'' below it, inside black dashed lines that bound a black truncated cone (sample about 400 µm wide); a much wider red truncated cone at the top with red dashed vertical lines down to a red double-headed arrow ``1,000 µm''.}
\figuredata{d}{figures/fig4d.csv}{figures/fig4de.py}{source data (tier 1): RP\_histMeasStar in YH\_MH\_traceabilityData\_fit\_mh\_v2.mat of the Supplementary Software, read by figures/fig4\_extract.py; tier-3 digitization for comparison in figures/fig4d\_digitized.csv}
\figuredata{e}{figures/fig4e.csv}{figures/fig4de.py}{source data (tier 1): RP\_histPredFinal in YH\_MH\_traceabilityData\_fit\_mh\_v2.mat, read by figures/fig4\_extract.py; tier-3 digitization for comparison in figures/fig4e\_digitized.csv}
\figuredata{f}{figures/fig4f.csv}{figures/fig4f.py}{source data (tier 1): RP\_priorx, RP\_priory in YH\_MH\_traceabilityData\_fit\_mh\_v2.mat, read by figures/fig4\_extract.py; tier-2 vector extraction for comparison in figures/fig4f\_vector.csv}
\figuredata{g}{figures/fig4g.csv}{figures/fig4g.py}{solid lines source data (tier 1): RP\_rp/RP\_p\_errN from the three .mat files of the Supplementary Software, read by figures/fig4\_extract.py; dashed lines vector paths in the PDF (tier 2; their source files are not in the Supplementary Software); full tier-2 extraction for comparison in figures/fig4g\_vector.csv}
\decompnote{Fig. 4g: the tick labelled ``1'' on the x axis is at the left end of the axis, one annotator-unit left of where the ticks labelled 10 and 20 put the value 1; the plotted solid lines start at the position of 1 annotator on the 10/20 scale. The legend does not match the data: by the authors' Supplementary Software, the red curve (legend ``e1088 (retina)'') is this paper's data set (the same fit as panels d–f), the green curve (legend ``This protocol (cortex)'') is the rOTO cortex comparison data set, and the blue curve (legend ``rOTO (cortex)'') is the e1088 retina data set; each solid curve matches its file within 0.4\,\% and differs from the other files by 34–190\,\%. See figures/fig4g.csv and the README.}

\noindent\textbf{Sample extraction and fixation.} Animals were anesthetized with isoflurane (Baxter) inhalation and perfused with 15\,ml cacodylate (Serva, Heidelberg, Germany) buffer (0.15\,M, pH 7.4) followed by 30\,ml fixative mixture containing 0.08\,M cacodylate (pH 7.4), 2.5\% paraformaldehyde (Sigma-Aldrich, St Louis, USA), 1.25\% glutaraldehyde (Serva) and 2\,mM calcium chloride (Sigma-Aldrich). The brain was removed from the skull with care to avoid mechanical irritation and post-fixed for 12 to 24\,h at 4\,°C in the fixative mixture. Then a biopsy punch (KAI Medical, Honolulu, USA) was used to extract samples sized about 1\,mm in diameter and about 2\,mm in length.

\textbf{Sample staining.} Large-volume \emph{en-bloc} staining was performed as follows (also see Supplementary Table 1 for details). Tissue was first immersed in 2\% \ce{OsO4} aqueous solution (Serva) buffered with cacodylate (0.15\,M, pH 7.4) at room temperature for 90\,min. The staining buffer was then replaced by 2.5\% ferrocyanide (Sigma-Aldrich) in 0.15\,M cacodylate buffer (pH 7.4) and incubated at room temperature for another 90\,min. Sequentially, the tissue was incubated in filtered thiocarbohydrazide (saturated aqueous solution at room temperature, Sigma-Aldrich) at 40\,°C for 45\,min, 2\% unbuffered \ce{OsO4} aqueous solution at room temperature for 90\,min and 1\% UA (Serva) aqueous solution at 4\,°C overnight. Double rinses in nanopure filtered water for 30\,min each were performed between the ferrocyanide and thiocarbohydrazide step, the thiocarbohydrazide and \ce{OsO4} step, and the \ce{OsO4} and UA step. On the next day, the tissue (still in UA solution) was warmed up to 50\,°C (oven) for 120\,min. After being washed twice in nanopure filtered water at room temperature for 30\,min, the tissue was incubated in a lead aspartate solution at 50\,°C for 120\,min. The lead aspartate solution was prepared by dissolving 0.066\,g lead nitrate (Sigma-Aldrich) in 10\,ml 0.03\,M aspartic acid (Serva) and pH adjusted to 5 with 1\,N KOH. The tissue was then washed twice in nanopure filtered water for 30\,min. The image in Supplementary Fig. 2b was taken from a tissue block stained with the same procedure described above except that the 120\,min 50\,°C incubation in UA was omitted.

\textbf{Sample embedding.} For embedding, samples were dehydrated through a graded ethanol series (50, 75, 100\%, 30\,min each, all cooled at 4\,°C) into pure acetone (3\,×\,100\%, 30\,min at room temperature) followed by epoxy-monomer infiltration by immersion into 1:1 mixtures of acetone and Spurr's resin (4.1\,g ERL 4221, 0.95\,g DER 736 and 5.9\,g NSA; Sigma-Aldrich) at room temperature for 12\,h or overnight on a rotator (samples were maintained in an open Eppendorf tube). Infiltrated samples were then incubated in pure resin with 1\% DMAE for 6\,h (in closed Eppendorf tubes) and placed in embedding moulds (Polyscience, Eppelheim, Germany) in a pre-warmed oven (70\,°C) for 48 to 72\,h.

\textbf{Sample screening.} Embedded samples were trimmed to a block face of 1.5\,×\,1.5\,mm\textsuperscript{2} and imaged in a scanning electron microscope with a field-emission cathode (QuantaFEG 200, FEI Company). An incident electron beam with a spot size of 2.9 and an energy of 2.8\,keV was used to scan across the samples at a pixel dwell time of 3.2\,µs and a pixel size of 6.1\,×\,6.1\,nm (Fig. 1b,c,e and Supplementary Fig. 2) or 12\,×\,12\,nm (Fig. 1d). Electrons were detected using a custom-designed detector based on a silicon diode (AXUV, International Radiation Detectors) combined with a custom-built current amplifier (courtesy of W. Denk, MPI Heidelberg), as described previously\textsuperscript{18,22}. To prevent charging for PATCO-stained samples, the chamber was kept at a pressure of 0.20\,mbar, while other samples were imaged at high vacuum. All samples were either mounted onto aluminium holders with electrical conductive glue (Henkel, Düsseldorf, Germany) or connected with conductive silver paint (SPI suppliers, West Chester, USA) to the metal stage.

High-resolution images (Fig. 2 and Supplementary Fig. 1) were taken from a tissue block mounted on an aluminium pin with electrically conductive glue. The tissue block was then trimmed at about 500-µm depth from the pia and at about 100\,µm from each other side, resulting in a block face with a size of about 970\,×\,890\,µm. The sample was then coated with a gold layer of 200-nm thickness, the block face was exposed again by carefully trimming off the surface with an ultra-microtome (Leica, UC7). EM micrographs were taken in a scanning electron microscope (Verios 460, FEI company) in high vacuum mode. An incident electron beam with an energy of 2.8\,keV and a current of 0.2\,nA was applied to scan across the sample at a pixel dwell time of 10\,µs with a pixel size of 1.35\,×\,1.35\,nm or 5.62\,×\,5.62\,nm.

\pdfpage{7}
\textbf{SBEM stack acquisition.} To test tissue cuttability and for performing tracing test experiments, we acquired a 3D SBEM data set (3\,×\,3 image tiles; single image size 2,048\,×\,1,768 pixels; lateral overlap between image tiles about 9\% in \emph{x}-direction and 6\% in \emph{y}-direction, respectively; 1,370 slices at nominal cutting thickness of 30\,nm; incident beam energy 2.8\,keV; dwell time 3.2\,µs; data set size 65\,×\,51\,×\,41\,µm\textsuperscript{3}; voxel size 12\,×\,12\,×\,30\,nm\textsuperscript{3}) from one of the samples processed with our protocol (Supplementary Table 1, Fig. 4a).

\textbf{Neurite reconstruction and traceability computations.} To generate tracing seed points, six image regions of size 2\,×\,2\,µm each were randomly selected from the data set. In these regions, all cross-sections of neuronal or glial processes were manually marked. Then, glial processes and cell bodies were excluded, leaving a total of 116 neurite processes. These were delivered to annotators as starting (seed) points. Annotators were asked to trace the respective neurite in all directions starting at the respective seed point, within a cubic bounding box of 10\,µm edge length, centred on the seed point. This yielded 116\,×\,7\,=\,812 single skeleton tracings, 22.7\,mm total path length. These were then analysed using the redundant skeleton consolidation procedure (RESCOP\textsuperscript{21}).

Briefly, each set of seven redundant skeletons was compared computing the number of `pro' votes (\emph{T}) and total votes (\emph{N}) for each edge in each skeleton tracing (parameters: `initial correspondence test radius' \emph{r}\textsubscript{p}: 225\,nm; r.m.s. node distance threshold Θ\,=\,375\,nm, see code in Supplementary Software). The resulting vote histogram was corrected for the redundancy (\emph{N}) of each tracing, yielding the measured vote histogram (Fig. 4d). Then, the underlying prior of edge probability \emph{p}(\emph{p}\textsubscript{e}) was determined by fitting the vote histogram under the simplifying assumption that tracing decisions were independent between tracers and locations (see ref. 21 for details). The best-fit prior and resulting predicted vote histogram are shown in Fig. 4e,f. Then, the predicted mean error-free path length as a function of the number of tracings per seed point (`tracing redundancy') was computed from the fitted prior \emph{p}(\emph{p}\textsubscript{e}) (solid lines in Fig. 4g) for the setting of annotators reconstructing an entire neuron/neurite; and for the setting of focused annotation where skeleton edges are repeatedly annotated until a set accuracy goal is reached (dashed lines in Fig. 4g, see ref. 21 for details). The code for this analysis (MATLAB, Mathworks) can be found in the Supplementary Software together with all tracings. For comparison, the published retina data\textsuperscript{21} was re-analysed based on the measured vote histogram using the same code, as was the comparison data set from rOTO-stained mouse cortex (Boergens \emph{et al.}, unpublished data set).
\decompnote{The Supplementary Software (RESCOP MATLAB code and three .mat files with the vote histograms, fits and error predictions; the tracings themselves are not in it) is in the source folder under esm\_zips/ and is the source of the Fig. 4d–g data here. The Supplementary Movie is not in the source folder.}

\textbf{Exploratory test experiments.} To ensure that the new protocol is readily applicable to other cortical areas and can be easily implemented by other experimentalists, we asked two students in the laboratory to use the protocol without direct supervision (just based on written instructions). Both were able to conduct successful stainings in somatosensory and parietal cortex (Supplementary Fig. 3 and Supplementary Table 2), though one sample showed charging in parts of the volume (Supplementary Fig. 3c). Successful protocol applications were required between 1 and about 10 training iterations for experienced and novice students, respectively.

\subsection*{References}
\begin{reflist}
1. Briggman, K. L. \& Bock, D. D. Volume electron microscopy for neuronal circuit reconstruction. \emph{Curr. Opin. Neurobiol.} \textbf{22,} 154–161 (2012).

2. Mikula, S. \& Denk, W. High-resolution whole-brain staining for electron microscopic circuit reconstruction. \emph{Nat. Methods} \textbf{12,} 541–546 (2015).\localcopy{mikula2015}

3. Helmstaedter, M. Cellular-resolution connectomics: challenges of dense neural circuit reconstruction. \emph{Nat. Methods} \textbf{10,} 501–507 (2013).

4. Palade, G. E. A study of fixation for electron microscopy. \emph{J. Exp. Med.} \textbf{95,} 285–298 (1952).

5. Watson, M. L. Staining of tissue sections for electron microscopy with heavy metals. \emph{J. Biophys. Biochem. Cytol.} \textbf{4,} 475–478 (1958).

6. Tapia, J. C. \emph{et al.} High-contrast en bloc staining of neuronal tissue for field emission scanning electron microscopy. \emph{Nat. Protoc.} \textbf{7,} 193–206 (2012).\localcopy{tapia2012}

7. Wai, J. S. M. \emph{et al.} A mechanistic insight leads to a greatly improved osmium-catalyzed asymmetric dihydroxylation process. \emph{J. Am. Chem. Soc.} \textbf{111,} 1123–1125 (1989).

8. Kwong, H.-L., Sorato, C., Ogino, Y., Chen, H. \& Sharpless, K. B. Preclusion of the "second cycle" in the osmium-catalyzed asymmetric dihydroxylation of olefins leads to a superior process. \emph{Tetrahedron Lett.} \textbf{31,} 2999–3002 (1990).

9. Kolb, H. C., VanNieuwenhze, M. S. \& Sharpless, K. B. Catalytic Asymmetric Dihydroxylation. \emph{Chem. Rev.} \textbf{94,} 2483–2547 (1994).

10. Locke, M., Krishnan, N. \& McMahon, J. T. A routine method for obtaining high contrast without staining sections. \emph{J. Cell Biol.} \textbf{50,} 540–544 (1971).

11. Khan, A. A., Riemersma, J. C. \& Booij, H. L. The reactions of osmium tetroxide with lipids and other compounds. \emph{J. Histochem. Cytochem.} \textbf{9,} 560–563 (1961).\localcopy{khan1961}

12. Korn, E. D. A chromatographic and spectrophotometric study of the products of the reaction of osmium tetroxide with unsaturated lipids. \emph{J. Cell Biol.} \textbf{34,} 627–638 (1967).\localcopy{korn1967}

13. White, D. L., Andrews, S. B., Faller, J. W. \& Barrnett, R. J. The chemical nature of osmium tetroxide fixation and staining of membranes by x-ray photoelectron spectroscopy. \emph{Biochim. Biophys. Acta} \textbf{436,} 577–592 (1976).\localcopy{white1976}

14. Subbaraman, L. R., Subbaraman, J. \& Behrman, E. J. \emph{Inorganic Chemistry} (Department of Biochemistry, Ohio State Univ.11, 2621–2627, 1972).
\decompnote{Reference 14 is printed in this garbled form (no article title; an affiliation in the place of the publisher; ``Univ.11''). It is presumably the 1972 \emph{Inorganic Chemistry} 11, 2621–2627 article by these authors.}\localcopy{subbaraman1972}

15. White, D. L., Mazurkiewicz, J. E. \& Barrnett, R. J. A chemical mechanism for tissue staining by osmium tetroxide-ferrocyanide mixtures. \emph{J. Histochem. Cytochem.} \textbf{27,} 1084–1091 (1979).\localcopy{white1979}

16. Seligman, A. M., Wasserkrug, H. L. \& Hanker, J. S. A new staining method (OTO) for enhancing contrast of lipid-{}-containing membranes and droplets in osmium tetroxide-{}-fixed tissue with osmiophilic thiocarbohydrazide(TCH). \emph{J. Cell Biol.} \textbf{30,} 424–432 (1966).\localcopy{seligman1966}

17. Willingham, M. C. \& Rutherford, A. V. The use of osmium-thiocarbohydrazide-osmium (OTO) and ferrocyanide-reduced osmium methods to enhance membrane contrast and preservation in cultured cells. \emph{J. Histochem. Cytochem.} \textbf{32,} 455–460 (1984).\localcopy{willingham1984}

18. Briggman, K. L., Helmstaedter, M. \& Denk, W. Wiring specificity in the direction-selectivity circuit of the retina. \emph{Nature} \textbf{471,} 183–188 (2011).

19. Holcomb, P. S. \emph{et al.} Synaptic inputs compete during rapid formation of the calyx of Held: a new model system for neural development. \emph{J. Neurosci.} \textbf{33,} 12954–12969 (2013).

20. Seligman, A. M., Hanker, J. S., Wasserkrug, H., Dmochowski, H. \& Katzoff, L. Histochemical demonstration of some oxidized macromolecules with thiocarbohydrazide (tch) or thiosemicarbazide (TSC) and osmium tetroxide. \emph{J. Histochem. Cytochem.} \textbf{13,} 629–639 (1965).\localcopy{seligman1965}

21. Helmstaedter, M., Briggman, K. L. \& Denk, W. High-accuracy neurite reconstruction for high-throughput neuroanatomy. \emph{Nat. Neurosci.} \textbf{14,} 1081–1088 (2011).

22. Helmstaedter, M. \emph{et al.} Connectomic reconstruction of the inner plexiform layer in the mouse retina. \emph{Nature} \textbf{500,} 168–174 (2013).
\end{reflist}

\subsection*{Acknowledgements}
We thank Kevin Briggman, Winfried Denk, Mark Ellisman, Kristen Harris and Shawn Mikula for fruitful discussions; Ali Karimi, Florian Drawitsch, Jakob Straehle and Yates Buckley for comments on the manuscript; Ali Karimi and Anjali Gour for performing protocol test experiments; Heiko Wissler for excellent help with tracer management and figure generation; Vera Nijveld for support with schematics; and Lisa Geiser, Raphael Jakoby, Raphael Kneissl, Vanessa Robl, Nathalie Schmidt and Victoria Schuhbeck for neurite reconstructions.

\subsection*{Author contributions}
Y.H. developed the protocol and carried out the experiments. P.L. contributed to protocol optimization and carried out experiments. Y.H. and M.H. analysed the data. Y.H. and M.H. wrote the paper.

\subsection*{Additional information}
\textbf{Supplementary Information} accompanies this paper at http://www.nature.com/naturecommunications

\noindent\textbf{Competing financial interests:} The authors declare no competing financial interests.

\noindent\textbf{Reprints and permission} information is available online at http://npg.nature.com/reprintsandpermissions/

\noindent\textbf{How to cite this article:} Hua, Y. \emph{et al.} Large-volume \emph{en-bloc} staining for electron microscopy-based connectomics. \emph{Nat. Commun.} 6:7923 doi: 10.1038/ncomms8923 (2015).

\medskip
\noindent This work is licensed under a Creative Commons Attribution 4.0 International License. The images or other third party material in this article are included in the article's Creative Commons license, unless indicated otherwise in the credit line; if the material is not included under the Creative Commons license, users will need to obtain permission from the license holder to reproduce the material. To view a copy of this license, visit http://creativecommons.org/licenses/by/4.0/
\decompnote{The licence statement is preceded by the Creative Commons ``cc by'' logo.}

\end{document}
